\documentclass[11pt]{article}
\usepackage[margin=0.85in]{geometry}
\usepackage{amsmath,booktabs,xurl,hyperref}
\usepackage[T1]{fontenc}
\setlength{\emergencystretch}{2em}
\title{Runtime Composition of the Longest CLF-Recovery Frame}
\author{Pan Dynamics Collision Avoidance Research}
\date{October 1, 2026}
\begin{document}
\maketitle

\section{Original measured frame}
The longest complete controller call in the fourteen-case campaign is
\textbf{5.539001 seconds}, in the 8-m/s \texttt{turningCrossing} scenario,
one-based frame 5, at simulation time 0.20 seconds. The driver starts its
clock immediately before calling the controller and stops it immediately
after return. This value includes input processing, formulation, all solver
calls and nonlinear admission; it excludes subsequent ODE45 integration,
offline geometry audit, plotting and export.

The saved trace establishes the following facts without reconstruction:
\begin{itemize}
\item The prediction has 212 holds at 50 ms: a total of 10.6 seconds.
Only eight holds (0.4 seconds) are the PCBF prefix; the remaining 204 holds
(10.2 seconds) are the optimized safety-completion tail.
\item The controller starts by restoring the shifted trajectory from the
changed measured ego state. Its exact RK4 successor does not equal the
independent ODE45 measurement, so the exact shifted-cache branch is unavailable.
\item Seven numerical optimizer calls are recorded. The one recorded SCvx
step uses three: a zero-CLF-slack SOCP, a relaxed-CLF SOCP, and a predictive-cost
SOCP fixing the minimizing first input. Its zero-slack solve reports infeasible
in the local model; the latter two solves return positive exit flags. The
remaining four calls must be correction QPs: those are the only counted
optimizer calls outside the recorded SCvx step in this source.
\item The primary safety solve is skipped because a zero-safety witness
already exists. Initial CLF slack is 8.166963; after the initial corrections,
the SCvx baseline slack is 6.006301. The admitted returned slack is 4.555671,
with zero hard residual and zero safety slack. CLF optimization is explicitly
incomplete and termination is \texttt{timeLimit}.
\end{itemize}
The five-second setting is a soft improvement budget. An in-flight solver,
rollout or correction is allowed to finish; the controller checks the clock
at selected boundaries. This explains how a call can exceed five seconds.
The original trace does not identify which operation accounts for its final
0.539001-second excess.

\section{Measured decomposition and its limit}
The original campaign did not save phase timers or the complete warm trajectory
before frame 5. Therefore its exact historical phase split cannot be recovered.
A diagnostic replay uses the recorded ego state, target epoch, previous issued
input and unchanged physical configuration. The first four controller calls
reconstruct a warm trajectory, always using the recorded physical states.
Budget-limited decisions differ from the original run: their largest input
component differences are 0.125000, 0.068771, 0.065142 and 0.107980. This is
\emph{not} a replay of the exact original optimizer state.

The resulting frame-5 input problem is frozen for three production replays
and three instrumented replays, using one single-thread MATLAB R2026a process.
Production calls take 5.040144, 5.494397 and 5.499686 seconds. Instrumented calls
take 5.029124, 5.013569 and 5.033861 seconds. All six have 212 holds, the same
issued input, zero hard residual and zero safety slack. Their issued input is
$[-0.010233318,-0.533951716]^T$, whereas the original is
$[-0.004935372,-0.469335934]^T$. Their final CLF slack is 5.27826.
The replay executes five SOCP calls over two SCvx steps and no correction QP
calls. This differs from the original seven-call path and precludes labeling
the following numbers as the exact decomposition of 5.539001 seconds.

The last warmed instrumented replay provides this additive breakdown:
\begin{center}
\begin{tabular}{lrr}
\toprule
Component & Seconds & Share \\
\midrule
SOCP numerical solves (five calls) & 2.337185 & 46.43\% \\
Affine problem construction and setup & 0.804729 & 15.99\% \\
Nonlinear rollout and safety evaluation & 0.860250 & 17.09\% \\
Nonlinear dynamic-defect retraction & 0.796426 & 15.82\% \\
Endpoint correction, excluding nested evaluations & 0.207087 & 4.11\% \\
First-step CLF correction and remaining orchestration & 0.028184 & 0.56\% \\
\midrule
Total & 5.033861 & 100.00\% \\
\bottomrule
\end{tabular}
\end{center}
The numerical-solve total consists of 0.695355 seconds for two infeasible
zero-slack attempts, 1.249449 seconds for two relaxed-CLF solves, and 0.392381
seconds for one fixed-first-input cost solve. The replay performs fourteen
nonlinear evaluations, two dynamic-defect retractions and five endpoint
correction invocations. The latter can return immediately; invocation count
is not a count of successful correction iterations.

Nested timers record inclusive and exclusive time, so a rollout called by an
endpoint correction is counted under rollout rather than counted twice.
Exclusive times sum to the root solve time within $10^{-9}$ seconds in all
three instrumented runs. The remaining outer controller time is assigned to
orchestration. Only an isolated source copy contains the timers and pass-through
solver wrappers. Production controller files are unchanged.

\section{Interpretation}
The expensive work is a long, repeatedly processed prediction. At this horizon,
the basic conic program has 1,918 variables and 1,278 dynamics equality rows;
the predictive-cost version has 2,130 variables and 214 cones. Measured replay
linear-inequality counts are approximately 2,975--2,980. The model includes
nonlinear combined-slip tires, yaw-dependent rectangle geometry, adaptive
interval separation and a terminal continuation. A single full rollout spans
2,120 RK4 integration intervals, before counting repeated evaluations and
variational dynamics.

In the measured replay, numerical optimization accounts for about 46\% of the
call. Formulation, nonlinear evaluation and correction account for almost all
of the rest. Improving only the SOCP engine would leave substantial cost in
repeated full-horizon work. The original frame additionally spends four QP
calls repairing trajectory feasibility, whose historical durations are unknown.
A zero PCBF slack does not imply that CLF work has finished: the original
returned CLF slack remains positive. Continuing to improve it is consistent
with the current objective hierarchy; the computational burden comes from the
size and repeated processing of the coupled continuation.

A local infeasible zero-slack SOCP is not a proof of globally impossible
nonlinear CLF dissipation. Neither the profile nor the original five-second
budget establishes a hard real-time deadline. Concurrent campaign load, the
reconstructed warm start and timer overhead all limit comparison of absolute
times. No algorithm change, new controller mode or road constraint was introduced
for this diagnostic.

The compact bundle \path{CLF_SLOW_FRAME_20261001/} contains the original frame,
reconstruction differences, replay checks, additive timings, individual events,
reproduction commands and source/raw hashes. Full MAT snapshots, generated
instrumentation and the derived PDF remain in the external research cache.
\end{document}
